\documentclass[12pt]{article}

% -------------------- PACKAGES --------------------
\usepackage[T1]{fontenc}
\usepackage{lmodern}
\usepackage{graphicx}
\usepackage{anysize}
\marginsize{2.1cm}{2.1cm}{1.8cm}{1.8cm}

\usepackage{url}
\usepackage{fancyhdr}
\usepackage{setspace}
\usepackage{lastpage}
\usepackage{titlesec}

% -------------------- HEADER / FOOTER --------------------
\pagestyle{fancy}
\fancyhf{}

\fancyhead[L]{\footnotesize UX Design, Game Design Internship}
\fancyhead[R]{\footnotesize \nouppercase{\leftmark}}

\fancyfoot[C]{\footnotesize 2025-26}
\fancyfoot[L]{\footnotesize Dept. of ECE, Vemana IT}
\fancyfoot[R]{Page \thepage\ of \pageref{LastPage}}

\renewcommand{\footrulewidth}{0.4pt}

% -------------------- SUBSECTION STYLE --------------------
\titleformat{\subsection}
{\bfseries\fontsize{14}{16}\selectfont}
{\thesubsection}{1em}{}

\onehalfspacing

% ==================== DOCUMENT ====================
\begin{document}

% ==================== CHAPTER 1 ====================

\section*{CHAPTER 1}
\markboth{INTRODUCTION}{}

% Fix numbering for this chapter
\setcounter{subsection}{0}
\renewcommand{\thesubsection}{1.\arabic{subsection}}

\begin{center}
\textbf{\Large INTRODUCTION}
\end{center}

"UX Design, Game Design Intern" was a paid internship course provided by Edutainer for the cost of Rs 6,000. This course was offered a clear view of Game Designing and UX Designing along with the Basic of computer science needed to understand the course easier.

This chapter fills in the details related to the course introduction, scope and its importance.

The acronyms computer science (CS), user experience (UX), and user interface (UI) are commonly used in the PDF going forward.

% -------------------- 1.1 --------------------
\subsection{Internship Scope and Objectives}

The course starts from the basics of computer science to ensure that the user/intern is well-knowledgeable regarding the details and make a smoother flow in the process of understanding the whole concept. 

The primary goal of the course being to get the intern to understand UX Design/Game Design. 

The UX Design covers the understanding of user behavior, developing user-friendly interfaces(UI elements) and developing intriguing user experiences (UX elements). 

The game design covers the understanding of concepts like mechanics, structure, and user engagement in game involving positions relating to UX research, interface design, and interaction design.

Design and prototyping are done with programs like Figma, Adobe XD, and Unity. User research, wireframing, prototyping, usability testing, and iterative design enhancements are some of the tasks.

% -------------------- 1.2 --------------------
\subsection{Relevance of UX and Game Design}

UX design is crucial to ensuring usability, accessibility, and user satisfaction in digital systems such as webpages, application software, etc. It directly affects the engagement and retention of system users.

Game design creates immersive experiences in interactive games by combining creative logic and storytelling to amuse the player.

Both domains are heavily utilised in real-world applications such as web platforms and webpages, mobile applications, gaming systems, and interactive technologies.

% -------------------- 1.3 --------------------
\subsection{Evolution of UX and Game Design Practices}

Basic UI design has been replaced by a user-centered, research-driven approach. While functionality was the primary focus of earlier systems, modern systems place a higher priority on usability and experience.

Simple 2D games gave way to intricate 3D settings with cutting-edge physics, graphics, and, more recently, artificial intelligence-powered interactions.

% -------------------- 1.4 --------------------
\subsection{Current Trends and Technologies}

Thanks to advancements like AI-driven programming and UI/UX customization, AR/VR technologies, and contemporary tools like Figma, Unity, and Unreal Engine, the field is changing.

The use of data-driven design techniques to enhance decision-making and user experience is growing.


% ==================== CHAPTER 2 ====================
\newpage

\section*{CHAPTER 2}
\markboth{ORGANISATION PROFILE}{}

\begin{center}
\textbf{\Large ORGANISATION PROFILE}
\end{center}

% Fix numbering for this chapter
\setcounter{subsection}{0}
\renewcommand{\thesubsection}{2.\arabic{subsection}}

In addition to answering questions like when and where the internship was completed, this chapter gives a general overview of the company and concentrates on its history, goals, services, and overall impact rather than highlighting individual experiences.

% -------------------- 2.1 --------------------
\subsection{Introduction}

Edutainer is a top supplier of online courses, career advancement programs, and virtual internships for undergraduates and working professionals. To help students succeed in the workplace. They specialize in job-related training, portfolio development, and internship projects. Microsoft Office, Python, Data Structures, UI/UX Design, React, and interview techniques are some of the subjects they typically cover. They combined custom learning solutions, such as LMS creation and API-based course delivery, in partnership with leading universities and training centers. They offer professional student counseling, job placement prep, and consulting assistance for skill development programs in addition to instruction.

\begin{figure}[htbp]
\centering

\includegraphics[width=7.5cm]{images/brand-logos/edutainer-logo.png}
\caption{Internship Provider - Edutainer Logo}
\end{figure}

% -------------------- 2.2 --------------------
\subsection{Vision and Mission of the Organization}

Edutainer envisions creating convenient, engaging and meaningful learning experiences through the integration of traditional educational methods with innovative technology.

Its mission is to establish an inclusive, unrestricted, and dynamic environment for learning.

% -------------------- 2.3 --------------------
\subsection{Programs and Services Offered}

Edutainer offers a wide range of programs across multiple domains, including UX Design, Game Design, Artificial Intelligence, and other emerging technologies.

The organization conducts virtual internships that simulate real-world working environments and provides structured skill development programs designed to enhance technical and practical competencies.

\begin{figure}[htbp]
\centering
\includegraphics[width=12.5cm]{images/graphics/Overviews/internship-overview.png}
\caption{Internship Overview}
\end{figure}

% -------------------- 2.4 --------------------
\subsection{Organizational Impact and Reach}

Edutainer achieved a major influence by addressing a large number of youngsters through its events and initiatives.

\begin{itemize}
\item 50,000+ students trained
\item 1000+ internships completed
\item High program completion rates
\item Multiple industry and academic partnerships
\end{itemize}

% -------------------- 2.5 --------------------
\subsection{Core Values}

The organization's approach to innovation and education is guided by a set of basic ideals.

\begin{itemize}
\item \textbf{Innovation} – Adopting new learning technologies and approaches.
\item \textbf{Accessibility} – Expanding the reach of education.
\item \textbf{Excellence} – Upholding high standards in content and delivery.
\item \textbf{Collaboration} – Creating alliances and learning communities.
\end{itemize}

% ==================== CHAPTER 3 ====================

\newpage
\section*{CHAPTER 3}
\markboth{WORK DONE / METHODOLOGY}{}

% Fix numbering for this chapter
\setcounter{subsection}{0}
\renewcommand{\thesubsection}{3.\arabic{subsection}}

\begin{center}
\textbf{\Large WORK DONE / METHODOLOGY}
\end{center}

The work completed during the internship and the technique used are presented in this chapter. It focuses on useful activities, equipment, and the methodical technique taken to do the assignment effectively.

The internship started on 2\textsuperscript{nd} February and lasted till May 31\textsuperscript{st}, 2026.

% -------------------- 3.1 --------------------
\subsection{Overview of Internship Work}

I focused on typical UX design and game design, structured learning, and technical documentation throughout my internship.

The tasks began with funadentals of cs work ,  in addition to arranging and updating instructional materials. Completing necessary classes, recording solutions, and keeping up with GitHub repositories were all part of this. One of the most important aspects of the internship was actively participating in quizzes, which involved reviewing questions and maintaining well-organised records for effective revision and clarity.

Another crucial component of the endeavour was content processing. This involved using AI-assisted techniques to enhance the information after using transcription technologies to convert audio and video data into text. The processed material was then carefully arranged into booklet, PDF, and note forms for better accessibility and understanding.

Assignments, notes, and official documents were all properly categorised according to the meticulous organization and upkeep of all resources via GitHub repositories. You may view the entire collection of work and structured material at: \url{https://github.com/Runarok/ux-game-intern/tree/main}.

\begin{figure}[htbp]
\centering
\includegraphics[width=12.5cm]{images/graphics/Overviews/workflow-cycle.png} 
\caption{Internship Workflow cycle}
\end{figure}

% -------------------- 3.2 --------------------
\subsection{Tools and Technologies Used}

A wide range of tools and technologies were used during the internship to support learning, design, and documentation processes.

For UX and design-related tasks, tools such as Figma and Canva were used for creating interfaces, layouts, and structured visual content. Unity was used for understanding basic game design concepts.

For content extraction and processing, transcription tools and speech-to-text software were utilized to convert video and audio materials into usable text formats.

AI tools and prompting interfaces were used to analyze, refine, and improve the generated content. Documentation and formatting were handled using PDF editors, note-taking applications, and booklet design tools.

GitHub was used as a central platform for organizing and managing all resources, including notes, assignments, and project files.

On the programming side, C language fundamentals were explored using GCC compiler and command-line environments, along with code editors.

\begin{figure}[htbp]
\centering
\includegraphics[width=15cm]{images/graphics/Overviews/tools-&-technologies.png} 
\caption{Tools and Technologies Used}
\end{figure}

% -------------------- 3.3 --------------------
\subsection{Design Process}

The internship followed a structured and iterative design and learning process.

The first stage involved understanding concepts, where core topics were introduced and studied in depth. This was followed by logical breakdown, where complex ideas were simplified into smaller, understandable components.

Next, application through examples was performed using pseudocode, flowcharts, and problem-solving scenarios. This helped in strengthening conceptual clarity.

In the UX context, the process included user research, wireframing, prototyping, and testing. Based on feedback and observations, iterative improvements were made to refine the designs.

Finally, documentation and reflection were carried out, where all learnings and solutions were organized into structured formats such as notes, PDFs, and GitHub repositories.

This structured workflow emphasizes concept understanding, logical breakdown, application, and documentation.

\begin{figure}[htbp]
\centering
\includegraphics[width=15cm]{images/graphics/Overviews/work-overview.png} 
\caption{Overview of Internship Work}
\end{figure}

% -------------------- 3.4 --------------------
\subsection{Game Design Analysis and Case Study}

This section focuses on understanding game design through the analysis of underlying systems rather than surface-level interface design. The emphasis is on how mechanics interact, how player decisions influence outcomes, and how meaningful gameplay emerges from simple rules.

A strong example of this approach is observed in \textit{Stone Story RPG}, which demonstrates how minimal presentation can support deeply engaging system design.

One of the most notable mechanics in the game is the \textbf{brewing (potion crafting) system}, which is built on combination logic and outcome-driven experimentation rather than explicit instruction. Instead of guiding the player step-by-step, the system encourages discovery through interaction.

Key characteristics of this system include:
\begin{itemize}
    \item \textbf{Combinational Design:} Different ingredients can be combined in varying sequences to produce distinct outcomes, creating a sense of experimentation and learning.
    \item \textbf{Emergent Gameplay:} The results of combinations are not always obvious, allowing players to discover patterns and optimize strategies over time.
    \item \textbf{Implicit Learning:} The system avoids heavy tutorials and instead teaches through feedback, repetition, and outcome observation.
    \item \textbf{Integration with Core Gameplay:} Brewing directly impacts combat efficiency, survival, and progression, making it a meaningful system rather than a side feature.
\end{itemize}

\begin{figure}[htbp]
\centering
\includegraphics[width=15cm]{images/stone-story-brewing.png}
\caption{Stone Story RPG – Brewing System and Combination Mechanics}
\end{figure}

This design highlights an important principle in game development: \textbf{depth through interaction rather than complexity through instruction}. The system remains simple at a glance but reveals depth as the player engages with it.

Additionally, exposure to idle and incremental games further reinforced understanding of long-term engagement systems. These games are built around continuous progression loops, automated resource generation, and exponential scaling, which sustain player interest over extended periods.

Such games typically emphasize:
\begin{itemize}
    \item Feedback loops and reward cycles
    \item Gradual unlocking of mechanics
    \item Optimization and efficiency improvements
\end{itemize}

A curated collection of personally recommended games, including Android titles and system-driven experiences worth exploring, is available at:
\url{https://github.com/Runarok/Guides/blob/main/Code%20Manual/Recommendations/Game%20Recommendations.pdf}

Overall, analyzing these systems provided a deeper understanding of how well-designed mechanics can create engaging experiences through player interaction, experimentation, and progression.

% -------------------- 3.5 --------------------
\subsection{Challenges Faced}

Several challenges were encountered during the internship.

One of the primary challenges was understanding complex concepts in computer science and UX design and breaking them down into simpler, structured formats.

Managing large amounts of content and organizing it efficiently was also difficult, especially while maintaining clarity and consistency.

Tool-related limitations, such as handling transcription accuracy and formatting outputs, required additional effort and refinement.

Another challenge was ensuring that the generated content was meaningful, accurate, and well-structured rather than just converted data.

% -------------------- 3.6 --------------------
\subsection{Solutions Implemented}

To overcome these challenges, a structured and methodical approach was adopted.

Complex concepts were broken down into smaller logical units, making them easier to understand and document.

AI-assisted tools were used effectively to refine and enhance content, ensuring clarity and accuracy. Iterative improvements were made after reviewing outputs multiple times.

A well-organized GitHub structure was maintained to manage all resources systematically, reducing confusion and improving accessibility.

Consistency in formatting and documentation was achieved by following a standard structure across all notes, assignments, and outputs.
\newpage

% ==================== CHAPTER 4 ====================

\section*{CHAPTER 4}
\markboth{RESULTS AND DISCUSSION}{}

% Fix numbering for this chapter
\setcounter{subsection}{0}
\renewcommand{\thesubsection}{4.\arabic{subsection}}

\begin{center}
\textbf{\Large RESULTS AND DISCUSSION}
\end{center}

The work that was evaluated and examined during the internship is included in this chapter. It highlights the successes, the skills acquired, and a critical assessment of the process as a whole.

\begin{figure}[htbp]
\centering
\includegraphics[width=15cm]{images/graphics/Overviews/results.png} 
\caption{Results}
\end{figure}

% -------------------- 4.1 --------------------
\subsection{Outcomes of the Work}

The internship's outcomes were arranged and recorded in a variety of ways.

A collection of educational materials, notes, and assignments were created and maintained using GitHub repositories. Content from a range of sources was refined, processed, and converted into readable, useful formats such as PDF documents and notes.

For UX design tasks, structured layouts and interface designs were produced using Figma, HTML coding and Canva. Basic understanding and application of game design concepts were also achieved through interaction-focused design approaches.

In general, the project improved the ability to present technical information and transform unstructured data into ordered knowledge.

% -------------------- 4.2 --------------------
\subsection{Learning Outcomes}

Both technical and analytical skills were greatly enhanced by the internship.

A solid grasp of UX thinking was attained by emphasizing usability, clarity, and user-centered design. The ability to produce well-organized and aesthetically pleasing outputs was improved by practical experience with design tools like Figma and Canva.

Constant practice of dissecting difficult ideas into more manageable logical steps helped to improve problem-solving abilities. Additionally, the capacity to analyze data, improve outputs, and methodically record outcomes was reinforced.

Additionally, the experience improved the ability to maintain organized workflows through self-learning, discipline, and consistency.

% -------------------- 4.3 --------------------
\subsection{Analysis}

The methodical approach used throughout the internship worked well to establish consistency and clarity. Learning was effectively reinforced by the process of comprehending concepts, dissecting them logically, applying them, and recording outcomes.

Efficiency and output quality were greatly increased by using tools like GitHub for organization and AI tools for content refinement.

However, certain limitations were observed. Early on, it was difficult to maintain consistency across documentation and manage enormous volumes of data. More work was required for some tools to yield accurate results, especially in the areas of formatting and content extraction.

To reduce the amount of time spent on repetitive tasks, workflow optimization has been improved. By focusing more on advanced design techniques and thoroughly examining game development tools, results could be further enhanced.

% ==================== CHAPTER 5 ====================

\newpage
\section*{CHAPTER 5}
\markboth{CONCLUSION}{}

% Fix numbering for this chapter
\setcounter{subsection}{0}
\renewcommand{\thesubsection}{5.\arabic{subsection}}

\begin{center}
\textbf{\Large CONCLUSION}
\end{center}

This chapter summarizes the overall internship experience, highlighting key outcomes, personal growth, and future directions.

\begin{figure}[htbp]
\centering
\includegraphics[width=15cm]{images/graphics/Overviews/Conclusion.png} 
\caption{Conclusion}
\end{figure}

% -------------------- 5.1 --------------------
\subsection{Summary}

The internship offered structured exposure to basic computer science concepts, game design, and UX design. It prioritised clarity, methodical learning, and logical reasoning over surface-level understanding.

The internship's goal of developing a solid foundation in the technical and analytical facets of the field was accomplished through a variety of tasks, including design work, content processing, documentation, and repository management.

% -------------------- 5.2 --------------------
\subsection{Personal Growth}

My skills and perspective were greatly enhanced by the internship.

Technical abilities such as UX thinking, organised problem-solving, and design tool usage were enhanced. Simultaneously, there was a discernible change in viewpoint—from merely finishing tasks to comprehending systems, procedures, and underlying logic.

Throughout the internship, one's capacity for independent learning, efficient information organization, and consistency all improved.

% -------------------- 5.3 --------------------
\subsection{Future Scope}

The fields of UX Design and Game Design are rapidly evolving with advancements in Artificial Intelligence, Augmented Reality, and immersive technologies.

There is significant scope for deeper exploration in user-centered design, interactive systems, and advanced game development. Future work can focus on building more complex and real-world applications using modern tools and frameworks.

This internship has created a strong base for pursuing further development and specialization in these domains.

\end{document}